\documentclass[compress]{beamer}
\usetheme{Montpellier}
\setbeamertemplate{navigation symbols}{horizontal}
%\useoutertheme{infolines}
\usecolortheme{dove}
%\usefonttheme{serif}

\usepackage[italian]{babel}
\usepackage{graphicx}
\usepackage{setspace}
\usepackage{tabularx}
\usepackage{booktabs}
\usepackage{fancyvrb}
\usepackage{xcolor}
\usepackage{multicol}
\usepackage{multirow}
\usepackage{tikz}
\usetikzlibrary{arrows.meta}
\usepackage{hyperref}

\definecolor{background}{RGB}{250, 250, 250}
\definecolor{back}{RGB}{247, 251, 255}
\definecolor{orangered2}{RGB}{238,64,0}
\definecolor{royalblue3}{RGB}{58,95,205}
 \setbeamerfont{subtitle}{size=\large, series=\bfseries}
\setbeamercolor{frametitle}{bg=background}
\setbeamertemplate{frametitle}{\color{template}\small \bfseries\insertframetitle\par\vskip-6pt\hrulefill}
\setbeamertemplate{navigation symbols}[default]
\definecolor{template}{RGB}{158, 158, 158}
\definecolor{latenti}{RGB}{54, 114, 89}
\definecolor{manifeste}{RGB}{179, 7, 27}
\definecolor{background}{RGB}{251, 251, 251}
\definecolor{highlight}{RGB}{18, 10, 143}
\definecolor{giallo}{RGB}{204 204 204}
\definecolor{blu}{RGB}{177 11 37}
\definecolor{orangered2}{RGB}{238,64,0}
\definecolor{typical}{RGB}{8, 69, 148}
\definecolor{springgreen}{RGB}{0,205,102}
\setbeamercolor{frametitle}{fg=blu, bg = white}
\setbeamercolor{item}{fg=giallo}
\setbeamercolor{caption name}{fg=blu}
\setbeamercolor{section name}{fg=white}
\setbeamercolor{subsection name}{fg=white}
\setbeamersize{text margin left=3mm,text margin right=3mm} 

\AtBeginSection[]{
  \begin{frame}{Outline}
    \tableofcontents[currentsection,currentsubsection]
  \end{frame}
}

\title[Experiment set up]{Misurazione Implicita in Psicologia:\\ Programmare gli esperimenti in Inquisit}
\date[PSQT, 2026]{Master di II Livello\\ Psicologia quantitativa: Misurazione, valutazione e analisi di variabili psicosociali}
\institute[]{11 settembre 2026, Padova}
\author[Ottavia M. Epifania]{\texorpdfstring{Ottavia M. Epifania, Ph.D.\newline\url{ottavia.epifania@unitn.it}}{Ottavia M. Epifania}}

\begin{document}

\begin{frame}[plain]
  \maketitle
\end{frame}

% =========================================================
\section*{Questo corso}
% =========================================================

\begin{frame}{Obiettivo della lezione}
Script Inquisit che contiene:

\vspace{3mm}
\begin{itemize}
  \item uno IAT funzionante;
  \item una scheda socio-demografica;
  \item alcune domande esplicite sullo stesso tema;
  \item due condizioni sperimentali basate sull'ordine di presentazione misure esplicite -- IAT.
\end{itemize}

\vspace{4mm}
\begin{exampleblock}{}
\centering
Gruppo 1: esplicite $\rightarrow$ IAT \qquad
Gruppo 2: IAT $\rightarrow$ esplicite
\end{exampleblock}
\end{frame}


% =========================================================
\section{Da un esperimento già pronto}
% =========================================================

\begin{frame}{Cosa ci serve}
\begin{itemize}
  \item Inquisit Lab: \href{https://www.millisecond.com/download/}{\textcolor{blue}{download}}.
  \item Age IAT dalla test library: \href{https://www.millisecond.com/download/library/iat/ageiat}{\textcolor{blue}{AGE IAT}}.
  \item Test library completa: \href{https://www.millisecond.com/download/library}{\textcolor{blue}{Millisecond Library}}.
  \item Documentazione: \href{https://www.millisecond.com/support/docs/current/html/main.htm}{\textcolor{blue}{Inquisit Help}}.
\end{itemize}

\vspace{2mm}
\begin{alertblock}{Attenzione}
La demo di Inquisit è valida per un mese. Dopo il mese di prova, si potrà ancora usare Inquisit per la programmazione degli esperimenti, ma il software non raccoglierà più i dati. 
\end{alertblock}
\end{frame}

\begin{frame}{Prima cosa: facciamolo correre}
\begin{enumerate}
  \item Aprite il file dell'Age IAT.
  \item Eseguite l'esperimento una volta.
  \item Guardate cosa vede un partecipante.
  \item Salvate il file dati con un nome univoco.
\end{enumerate}

\vspace{4mm}
\begin{exampleblock}{Esempio}
\centering
\texttt{dati-Mario-Rossi.dat}
\end{exampleblock}

\vspace{3mm}
Cosa cambiereste prima di usare davvero questo esperimento?
\end{frame}

% =========================================================
\section{La logica di Inquisit}
% =========================================================

\begin{frame}{Il workflow di Inquisit}
\centering
\resizebox{!}{0.67\textheight}{%
\begin{tikzpicture}[
  nodebox/.style={draw, rounded corners=4pt, minimum width=4.9cm, minimum height=0.95cm, align=center, fill=gray!8},
  mainbox/.style={draw, rounded corners=4pt, minimum width=4.9cm, minimum height=0.95cm, align=center, fill=template!16},
  finalbox/.style={draw, rounded corners=4pt, minimum width=4.9cm, minimum height=0.95cm, align=center, fill=orange!20},
  flowarrow/.style={-{Latex[length=2.2mm]}, thick, draw=gray!70}
]
\node[nodebox] at (0,0) (objects) {\textbf{OGGETTI}\\[-1mm]\scriptsize Cosa voglio mostrare?\\[-1mm]\scriptsize \texttt{<item></item>}, \texttt{<text></text>}, \texttt{<picture></picture>}};
\node[nodebox] at (0,-1.35) (trial) {\textbf{TRIAL}\\[-1mm]\scriptsize Cosa succede quando un oggetto è modificato?\\[-1mm]\scriptsize \texttt{<trial></trial>}};
\node[mainbox] at (0,-2.70) (block) {\textbf{BLOCK}\\[-1mm]\scriptsize Come organizzo in modo coerente i trial?\\[-1mm]\scriptsize \texttt{<block></block>}};
\node[finalbox] at (0,-4.05) (expt) {\textbf{EXPERIMENT}\\[-1mm]\scriptsize In which order are blocks presented?\\[-1mm]\scriptsize \texttt{<expt></expt>}};
\draw[flowarrow] (objects.south) -- (trial.north);
\draw[flowarrow] (trial.south) -- (block.north);
\draw[flowarrow] (block.south) -- (expt.north);
\end{tikzpicture}}
\end{frame}

\begin{frame}{Cosa faremo}
\centering
\resizebox{0.95\textwidth}{!}{%
\begin{tikzpicture}[
	phase/.style={
		draw,
		rounded corners=4pt,
		minimum width=3.4cm,
		minimum height=1.15cm,
		align=center,
		fill=gray!8
	},
	phasehi/.style={
		draw,
		rounded corners=4pt,
		minimum width=3.4cm,
		minimum height=1.15cm,
		align=center,
		fill=template!18
	},
	arrow/.style={
		-{Latex[length=2.2mm]},
		thick,
		draw=gray!70
	}
	]
	
	\node[phase] at (0,0.5) (one) {
		\textbf{1. Modifica}\\[-1mm]
		\scriptsize Punto di partenza: Uno IAT esistente
	};
	
	\node[phase] at (4.0,-1.5) (two) {
		\textbf{2. Costruisci}\\[-1mm]
		\scriptsize aggiungi le misure esplicite
	};
	
	\node[phasehi] at (8,-3.5) (three) {
		\textbf{3. Organizza}\\[-1mm]
		\scriptsize Ordine di presentazione controbilanciato
	};
	
	\draw[arrow] (one.south) |- (two.west);
	\draw[arrow] (two.south) |- (three.west);
	
	\end{tikzpicture}}

\end{frame}

% =========================================================
\section{Modificare uno IAT esistente}
% =========================================================

\begin{frame}[fragile]{Objects: etichette e stimoli}
\begin{columns}[T]
\begin{column}{.50\linewidth}
\begin{small}
\begin{Verbatim}[commandchars=\\\{\}]
<item \textcolor{red}{attributeAlabel}>
/1 = "Good"
</item>

<item \textcolor{red}{attributeA}>
/1 = "Joy"
/2 = "Love"
/3 = "Peace"
/4 = "Wonderful"
/5 = "Pleasure"
/6 = "Glorious"
/7 = "Laughter"
/8 = "Happy"
</item>
\end{Verbatim}
\end{small}
\end{column}
\begin{column}{.48\linewidth}
Qui diciamo a Inquisit \textbf{che cosa esiste}: categorie, parole, immagini, istruzioni.

\vspace{4mm}
Per aggiungere uno stimolo basta aggiungere una riga:
\begin{Verbatim}
/9 = "Bello"
\end{Verbatim}
\end{column}
\end{columns}
\end{frame}

\begin{frame}[fragile]{Objects: stimoli target}
\begin{columns}[T]
\begin{column}{.43\linewidth}
\begin{small}
\begin{Verbatim}[commandchars=\\\{\}]
<item \textcolor{red}{targetAlabel}>
/1 = "Young"
</item>

<item \textcolor{red}{targetA}>
/1 = "yf1.jpg"
/2 = "yf4.jpg"
/3 = "yf5.jpg"
/4 = "ym2.jpg"
/5 = "ym3.jpg"
/6 = "ym5.jpg"
</item>
\end{Verbatim}
\end{small}
\end{column}
\begin{column}{.55\linewidth}
Le immagini devono essere nella stessa cartella dello script.

\vspace{4mm}
Anche qui la logica è la stessa: una lista di elementi che verranno poi richiamati da \texttt{<picture>} e da \texttt{<trial>}.
\end{column}
\end{columns}
\end{frame}

\begin{frame}[fragile]{Appearance: testo e immagini}
\begin{columns}[T]
\begin{column}{.50\linewidth}
\begin{Verbatim}[commandchars=\\\{\}]
<text \textcolor{red}{attributeA}>
/ items = \textcolor{red}{attributeA}
/ fontstyle = ("Arial", 5%)
/ txcolor = green
</text>
\end{Verbatim}
\end{column}
\begin{column}{.50\linewidth}
\begin{Verbatim}[commandchars=\\\{\}]
<picture \textcolor{red}{targetA}>
/ items = \textcolor{red}{targetA}
/ size = (20%, 20%)
</picture>
\end{Verbatim}
\end{column}
\end{columns}

\vspace{4mm}
\centering
\texttt{<item>} definisce il contenuto; \texttt{<text>} e \texttt{<picture>} definiscono come appare.
\end{frame}

\begin{frame}[fragile]{Behavior: il trial}
\begin{small}
\begin{Verbatim}[commandchars=\\\{\}]
<trial \textcolor{red}{attributeA}>
/ validresponse = ("E", "I")
/ correctresponse = ("E")
/ stimulusframes = [1 = \textcolor{red}{attributeA}, \textcolor{blue}{errorReminder}]
/ posttrialpause = parameters.ISI
</trial>
\end{Verbatim}
\end{small}

\vspace{3mm}
Il trial definisce il comportamento: quali risposte sono valide, qual è la risposta corretta, quali stimoli vengono mostrati e cosa succede dopo.
\end{frame}

\begin{frame}[fragile]{Blocks: mettere insieme i trial}
\begin{small}
\begin{Verbatim}[commandchars=\\\{\}]
<block \textcolor{red}{attributepractice}>
/ bgstim = (\textcolor{red}{attributeAleft}, \textcolor{red}{attributeBright})
/ trials = [
  1 = instructions;
  2-21 = random(\textcolor{blue}{attributeA}, \textcolor{blue}{attributeB});
]
/ errormessage = true(error,200)
/ responsemode = correct
</block>
\end{Verbatim}
\end{small}

\vspace{2mm}
Il blocco decide quanti trial fare e in che ordine presentarli.
\end{frame}

\begin{frame}{Esercizio 1: riconoscere la struttura}
Nel file Age IAT:

\vspace{2mm}
\begin{enumerate}
  \item traducete le categorie che vede il partecipante;
  \item aggiungete una parola per ciascuna categoria attributo;
  \item aggiungete una immagine per ciascuna categoria target;
  \item cambiate dimensione/font degli stimoli;
  \item cambiate i tasti di risposta;
  \item modificate il numero di trial in un blocco di pratica.
\end{enumerate}

\vspace{3mm}
Per ogni modifica chiedetevi: sto toccando \texttt{<item>}, \texttt{<text>/<picture>}, \texttt{<trial>} o \texttt{<block>}?
\end{frame}

\begin{frame}[fragile]{Mini-task: togliere il reminder di errore}
\begin{small}
\begin{Verbatim}[commandchars=\\\{\}]
<trial attributeA>
/ validresponse = ("E", "I")
/ correctresponse = ("E")
/ stimulusframes = [1 = attributeA, \textcolor{red}{errorReminder}]
/ posttrialpause = parameters.ISI
</trial>
\end{Verbatim}
\end{small}

\vspace{2mm}
Per toglierlo:
\begin{small}
\begin{Verbatim}
/ stimulusframes = [1 = attributeA]
\end{Verbatim}
\end{small}

\vspace{2mm}
Questa modifica va ripetuta nei trial in cui il reminder è presente.
\end{frame}

\begin{frame}[fragile]{Mini-task: aggiungere una schermata reminder}
\centering
\resizebox{0.78\textwidth}{!}{%
\begin{tikzpicture}[
  b/.style={draw, rounded corners=4pt, minimum width=3.7cm, minimum height=.8cm, align=center, fill=gray!8},
  a/.style={-{Latex[length=2mm]}, thick, draw=gray!70}
]
\node[b] at (0,0) (obj) {\textbf{1. OBJECT}\\[-1mm]\scriptsize \texttt{<text reminder>}};
\node[b] at (4.2,0) (tri) {\textbf{2. BEHAVIOR}\\[-1mm]\scriptsize \texttt{<trial reminder>}};
\node[b] at (8.4,0) (blo) {\textbf{3. POSITION}\\[-1mm]\scriptsize add to \texttt{<block>}};
\draw[a] (obj.east) -- (tri.west);
\draw[a] (tri.east) -- (blo.west);
\end{tikzpicture}}

\vspace{5mm}
\begin{small}
\begin{Verbatim}[commandchars=\\\{\}]
<text reminder>
/ items = ("Controlla le categorie - premi SPAZIO per cominciare")
/ position = (50%, 75%)
/ color = (255, 255, 255)
/ fontstyle = ("Arial", 3%, false, false, false, false, 5, 1)
</text>
\end{Verbatim}
\end{small}
\end{frame}

\begin{frame}[fragile]{Reminder: trial e block}
\begin{columns}[T]
\begin{column}{.48\linewidth}
\begin{small}
\begin{Verbatim}[commandchars=\\\{\}]
<trial \textcolor{blue}{reminder}>
/ correctresponse = (" ")
/ frames = [1=reminder]
/ responsemode = correct
</trial>
\end{Verbatim}
\end{small}
\end{column}
\begin{column}{.50\linewidth}
\begin{small}
\begin{Verbatim}[commandchars=\\\{\}]
<block attributepractice>
/ trials = [
  1=instructions;
  \textcolor{red}{2=reminder;}
  3-22 = random(attributeA, attributeB);
]
</block>
\end{Verbatim}
\end{small}
\end{column}
\end{columns}

\vspace{3mm}
\begin{alertblock}{Attenzione}
Quando aggiungete un trial, aggiornate anche la numerazione dei trial successivi.
\end{alertblock}
\end{frame}

% =========================================================
\section{Aggiungere le misure esplicite}
% =========================================================

\begin{frame}{Costruire una nuova scheda socio demografica}


\vspace{2mm}
\centering
\resizebox{!}{0.46\textheight}{%
\begin{tikzpicture}[
  bx/.style={draw, rounded corners=4pt, minimum width=4cm, minimum height=.85cm, align=center, fill=template!13},
  ar/.style={-{Latex[length=2mm]}, thick, draw=gray!70}
]
\node[bx] at (0,0) (q) {\textbf{QUESTIONS}\\[-1mm]\scriptsize \texttt{radiobuttons}, \texttt{dropdown}, \texttt{slider}};
\node[bx] at (0,-1.2) (p) {\textbf{SURVEY PAGE}\\[-1mm]\scriptsize \texttt{<surveypage></surveypage>}};
\node[bx] at (0,-2.4) (b) {\textbf{BLOCK}\\[-1mm]\scriptsize \texttt{<block></block>}};
\node[bx] at (0,-3.6) (e) {\textbf{EXPERIMENT}\\[-1mm]\scriptsize add to \texttt{<expt></expt>}};
\draw[ar] (q.south) -- (p.north);
\draw[ar] (p.south) -- (b.north);
\draw[ar] (b.south) -- (e.north);
\end{tikzpicture}}
\end{frame}

\begin{frame}{Domande esplicite}
Le opzioni più comuni per creare domande esplicite sono:

\vspace{2mm}
\begin{itemize}
  \item \texttt{checkboxes}: più risposte possibili;
  \item \texttt{dropdown}: una risposta da menu a tendina;
  \item \texttt{radiobuttons}: una risposta tra opzioni visibili;
  \item \texttt{slider}: risposta su scala continua o ordinale.
\end{itemize}

\vspace{3mm}
Le domande possono essere messe sulla stessa \texttt{surveypage} o divise su pagine diverse.
\end{frame}

\begin{frame}{Cosa costruiamo}
\textbf{Scheda socio-demografica}
\begin{itemize}
  \item genere;
  \item età;
  \item occupazione.
\end{itemize}

\vspace{3mm}
\textbf{Domande esplicite}
\begin{itemize}
  \item simpatia verso le persone giovani;
  \item simpatia verso le persone anziane.
\end{itemize}

\vspace{3mm}
Scheda e domande esplicite andranno su due pagine diverse.
\end{frame}

\begin{frame}[fragile]{Genere, età, occupazione}
\begin{small}
\begin{Verbatim}[commandchars=\\\{\}]
<radiobuttons \textcolor{blue}{genere}>
/ caption = "Genere"
/ options = ("Femmina", "Maschio", "Altro")
/orientation = horizontal
</radiobuttons>
\end{Verbatim}
\end{small}

\begin{small}
\begin{Verbatim}[commandchars=\\\{\}]
<dropdown \textcolor{blue}{eta}>
/ caption = "Età in anni compiuti"
/ options = ("17", "18", "19", "20", ..., "più di 45")
/ optionvalues = ("17", "18", "19", "20", ..., "più di 45")
</dropdown>
\end{Verbatim}
\end{small}

\begin{small}
\begin{Verbatim}[commandchars=\\\{\}]
<checkboxes \textcolor{blue}{occupazione}>
/ caption = "Scegli la tua occupazione"
/ options = ("Lavoratore", "Studente", "Disoccupato", "Dottorando", 
"Assegnista")
/ optionvalues = ("lav", "stud", "dis", "dott", "ass")
</checkboxes>
\end{Verbatim}
\end{small}
\end{frame}

\begin{frame}[fragile]{La pagina socio-demografica}
\begin{small}
\begin{Verbatim}[commandchars=\\\{\}]
<surveypage \textcolor{blue}{demografica}>
/ questions = [1=\textcolor{red}{genere}; 2=\textcolor{red}{eta}; 3=\textcolor{red}{occupazione}]
/ txcolor = black
/ finishlabel = "Pagina Successiva"
/ nextbuttonposition = (80,90)
/ showpagenumbers = false
/ showquestionnumbers = false
</surveypage>

<block \textcolor{blue}{demografica}>
/ trials = [1=\textcolor{red}{demografica}]
/ screencolor = (170, 170, 170)
</block>
\end{Verbatim}
\end{small}
\end{frame}

\begin{frame}[fragile]{Valutazione esplicita}
\begin{footnotesize}
\begin{Verbatim}[commandchars=\\\{\}]
<slider \textcolor{blue}{giovani}>
/ caption = "Le persone giovani mi stanno simpatiche"
/ labels = ("1~nCompletamente in disaccordo", "2", "3", 
"4~nCompletamente d'accordo")
/ range = (1, 4)
/ defaultresponse = "1"
/ increment = 0.5
/ position = (5%, 40%)
/ showticks = true
/ slidersize = (70%, 10%)
/ txcolor = black
/ fontstyle = ("Arial", 20)
/ responsefontstyle = ("Arial",16)
</slider>
\end{Verbatim}
\end{footnotesize}

\vspace{2mm}
Stessa struttura per \texttt{anziani}, cambiando etichetta e testo della domanda.
\end{frame}

\begin{frame}[fragile]{La pagina delle domande esplicite}
\begin{small}
\begin{Verbatim}[commandchars=\\\{\}]
<surveypage \textcolor{blue}{simpatia}>
/ questions = [1=\textcolor{blue}{giovani}; 2=\textcolor{blue}{anziani}]
/ txcolor = black
/ finishlabel = "Pagina Successiva"
/ nextbuttonposition = (80,90)
/ showpagenumbers = false
/ showquestionnumbers = false
</surveypage>

<block \textcolor{blue}{simpatia}>
/ trials = [1=\textcolor{red}{simpatia}]
/ screencolor = (170, 170, 170)
</block>
\end{Verbatim}
\end{small}
\end{frame}

\begin{frame}{Esercizio 2: costruire le misure esplicite}
Create:

\vspace{2mm}
\begin{enumerate}
  \item una \texttt{surveypage} socio-demografica con genere, età e titolo di studio;
  \item una \texttt{surveypage} con due domande esplicite sul tema dello IAT;
  \item i rispettivi \texttt{<block>}.
\end{enumerate}

\vspace{4mm}
Non devono ancora comparire nell'esperimento: li assembleremo nel prossimo passaggio.
\end{frame}

% =========================================================
\section{Assemblare l'esperimento}
% =========================================================

\begin{frame}{Abbiamo tutti i pezzi}
\centering
\resizebox{0.85\textwidth}{!}{%
\begin{tikzpicture}[
  comp/.style={draw, rounded corners=4pt, minimum width=3.3cm, minimum height=1cm, align=center, fill=gray!8},
  ex/.style={draw, rounded corners=4pt, minimum width=3.7cm, minimum height=1cm, align=center, fill=orange!20},
  ar/.style={-{Latex[length=2mm]}, thick, draw=gray!70}
]
\node[comp] at (0,0) (iat) {\textbf{IAT}\\[-1mm]\scriptsize practice + test blocks};
\node[comp] at (4.5,0) (exp) {\textbf{EXPLICIT}\\[-1mm]\scriptsize demografica + simpatia};
\node[ex] at (2.25,-2.0) (e) {\textbf{\texttt{<expt>}}\\[-1mm]\scriptsize experiment order};
\draw[ar] (iat.south) -- (e.north west);
\draw[ar] (exp.south) -- (e.north east);
\end{tikzpicture}}

\vspace{5mm}
\textbf{in quale ordine presento i blocchi?}
\end{frame}

\begin{frame}{La manipolazione finale}
\centering
\resizebox{0.90\textwidth}{!}{%
\begin{tikzpicture}[
  root/.style={draw, rounded corners=4pt, minimum width=3.6cm, minimum height=.8cm, align=center, fill=template!15},
  group/.style={draw, rounded corners=4pt, minimum width=3.9cm, minimum height=.8cm, align=center, fill=orange!18},
  block/.style={draw, rounded corners=4pt, minimum width=3.9cm, minimum height=.8cm, align=center, fill=gray!8},
  ar/.style={-{Latex[length=2mm]}, thick, draw=gray!70}
]
\node[root] at (0,0) (subj) {PARTICIPANT};
\node[root] at (0,-1.05) (groups) {\texttt{/ groups}};
\node[group] at (-2.5,-2.25) (g1) {\textbf{Group 1}\\[-1mm]\scriptsize \texttt{(1 of 2)}};
\node[group] at (2.5,-2.25) (g2) {\textbf{Group 2}\\[-1mm]\scriptsize \texttt{(2 of 2)}};
\node[block] at (-2.5,-3.55) (g1a) {Explicit measures};
\node[block] at (-2.5,-4.65) (g1b) {IAT};
\node[block] at (2.5,-3.55) (g2a) {IAT};
\node[block] at (2.5,-4.65) (g2b) {Explicit measures};
\draw[ar] (subj.south) -- (groups.north);
\draw[ar] (groups.south) -- (g1.north);
\draw[ar] (groups.south) -- (g2.north);
\draw[ar] (g1.south) -- (g1a.north);
\draw[ar] (g1a.south) -- (g1b.north);
\draw[ar] (g2.south) -- (g2a.north);
\draw[ar] (g2a.south) -- (g2b.north);
\end{tikzpicture}}
\end{frame}

\begin{frame}[fragile]{Prima: lo IAT originale usa già due gruppi}
\begin{columns}[T]
\begin{column}{.49\linewidth}
\begin{footnotesize}
\begin{Verbatim}[commandchars=\\\{\}]
<expt>
/ onexptbegin = [
 values.conditionOrder = "c-ic"
]
/ groups = \textcolor{blue}{(1 of 2)}
/ blocks = [
 1=targetcompatiblepractice;
 2=attributepractice;
 3=compatibletest1;
 4=compatibletest2;
 5=targetincompatiblepractice;
 6=incompatibletest1;
 7=incompatibletest2;
 8=summary; 9=end;
]
</expt>
\end{Verbatim}
\end{footnotesize}
\end{column}
\begin{column}{.49\linewidth}
\begin{footnotesize}
\begin{Verbatim}[commandchars=\\\{\}]
<expt>
/ onexptbegin = [
 values.conditionOrder = "ic-c"
]
/ groups = \textcolor{blue}{(2 of 2)}
/ blocks = [
 1=targetincompatiblepractice;
 2=attributepractice;
 3=incompatibletest1;
 4=incompatibletest2;
 5=targetcompatiblepractice;
 6=compatibletest1;
 7=compatibletest2;
 8=summary; 9=end;
]
</expt>
\end{Verbatim}
\end{footnotesize}
\end{column}
\end{columns}
\end{frame}

\begin{frame}{Scelta didattica}
Per l'esercitazione di oggi fissiamo l'ordine interno dello IAT.

\vspace{4mm}
Non manipoliamo più anche C--IC vs IC--C

\vspace{4mm}
\begin{exampleblock}{Manipoliamo una sola cosa}
\centering
l'ordine tra misure esplicite e IAT
\end{exampleblock}
\end{frame}

\begin{frame}[fragile]{Group 1: esplicite prima dello IAT}
\begin{footnotesize}
\begin{Verbatim}[commandchars=\\\{\}]
<expt>
/ onexptbegin = [
 values.conditionOrder = "c-ic"
]
/ preinstructions = (iatintro)
\textcolor{red}{/ groups = (1 of 2)}
/ blocks = [
 \textcolor{blue}{1=demo_istruzioni; 2=demografica; 3=simpatia;}
 4=IAT1; 5=targetcompatiblepractice;
 6=IAT2; 7=attributepractice;
 8=IAT3; 9=compatibletest1;
 10=IAT_pause; 11=compatibletest2;
 12=IAT4; 13=targetincompatiblepractice;
 14=IAT5; 15=incompatibletest1;
 16=IAT_pause; 17=incompatibletest2;
 18=summary; 19=end;
]
</expt>
\end{Verbatim}
\end{footnotesize}
\end{frame}

\begin{frame}[fragile]{Group 2: esplicite dopo lo IAT}
\begin{footnotesize}
\begin{Verbatim}[commandchars=\\\{\}]
<expt>
/ onexptbegin = [
 values.conditionOrder = "c-ic"
]
/ preinstructions = (iatintro)
\textcolor{red}{/ groups = (2 of 2)}
/ blocks = [
 1=IAT1; 2=targetcompatiblepractice;
 3=IAT2; 4=attributepractice;
 5=IAT3; 6=compatibletest1;
 7=IAT_pause; 8=compatibletest2;
 9=IAT4; 10=targetincompatiblepractice;
 11=IAT5; 12=incompatibletest1;
 13=IAT_pause; 14=incompatibletest2;
 \textcolor{blue}{15=demo_istruzioni; 16=demografica; 17=simpatia;}
 18=summary; 19=end;
]
</expt>
\end{Verbatim}
\end{footnotesize}
\end{frame}

\begin{frame}{Esercizio finale in aula}
Avete uno script quasi completo.

\vspace{3mm}
Il vostro compito è assemblarlo in due condizioni:

\vspace{3mm}
\begin{center}
\begin{tabular}{lll}
\toprule
\textbf{Gruppo} & \textbf{Ordine} & \textbf{Cosa modificare} \\
\midrule
1 & Esplicite $\rightarrow$ IAT & \texttt{/ groups} e \texttt{/ blocks} \\
2 & IAT $\rightarrow$ Esplicite & \texttt{/ groups} e \texttt{/ blocks} \\
\bottomrule
\end{tabular}
\end{center}

\vspace{3mm}
Quello che conta non è copiare righe: è capire dove si decide l'ordine sperimentale.
\end{frame}

% =========================================================
\section{Dati e analisi}
% =========================================================

\begin{frame}[fragile]{Salvare i dati}
\begin{small}
\begin{Verbatim}[commandchars=\\\{\}]
<data>
/ columns = (build, computer.platform, date, time,
 subject, group, session,
 blockcode, blocknum, trialcode, trialnum,
 values.conditionOrder,
 response, correct, latency,
 stimulusnumber, stimulusitem,
 expressions.da, expressions.db, expressions.d,
 expressions.percentcorrect,
 expressions.propRT300, expressions.excludeCriteriaMet)
</data>
\end{Verbatim}
\end{small}

\vspace{2mm}
Per analizzare i dati ci servono soprattutto: partecipante, gruppo, blocco, trial, risposta, correttezza e latenza.
\end{frame}

\begin{frame}{Summary dei risultati}
Inquisit può mostrare un feedback immediato con il \emph{D-score}.

\vspace{3mm}
\begin{exampleblock}{Pro}
Feedback immediato e possibilità di salvare il punteggio nel file dati.
\end{exampleblock}

\begin{alertblock}{Contro}
Da evitare prima di eventuali questionari espliciti. Inoltre, per analisi replicabili, è meglio ricalcolare il \emph{D-score} a partire dai dati grezzi.
\end{alertblock}
\end{frame}

% =========================================================
\section{A casa}
% =========================================================

\begin{frame}{Consegna}
Create uno IAT per investigare un tema di vostro interesse.

\vspace{3mm}
Aggiungete:
\begin{itemize}
  \item almeno 4 domande socio-demografiche;
  \item almeno 3 domande esplicite sullo stesso tema dello IAT;
  \item due condizioni sperimentali: esplicite prima vs esplicite dopo;
  \item esportazione dei dati trial-by-trial.
\end{itemize}

\end{frame}

\end{document}
